\documentclass[12pt]{article}
\usepackage[margin=1in]{geometry}
\usepackage[authoryear,round]{natbib}

\title{Referee Report for Oehmke and Opp (2025)}
\author{Siddhartha Lewis-Hayre}
\date{\today}

\begin{document}
\maketitle

% Structure follows Berk, Harvey, and Hirshleifer,
% Preparing a Referee Report: Guidelines and Perspectives.
% https://www.aeaweb.org/content/file?id=222

\section{Summary}

How should socially responsible (SR) funds invest to make an impact? There has been a large rise in socially responsible investing; however, critics are skeptical of its impact. They argue that divestment will not distort allocations as a return-maximizing fund will simply replace the divested capital. As such, socially responsible investment will merely generate worse returns without any social impact.

\citet{oehmke_opp_2025} acknowledge these critiques of SR funds with a model that investigates whether these funds have value even if return-maximizing funds have deep pockets and socially responsible investments have lower returns. With this setup, they show SR funds that try to maximize some combination of returns and the social good of their portfolio companies (what they call a narrow mandate) have no impact on the allocation. However, despite their setup, which may seem hopeless for SR funds, they show that SR funds that try to maximize some combination of returns and social good for all companies (what they call an impact mandate) do have a positive effect on the allocation. This positive impact can be achieved by investing in companies that previously used ``dirty" technologies (ones that create negative externalities) and requiring them to adopt governance rules mandating they use some ``clean'' technologies (ones that do not create externalities).

Finally, the paper considers what happens when you include financial constraints \`a la \citet{holmstrom_tirole_1997}. Consider a firm that would be financially constrained if using the clean technology. If an SR fund with an impact mandate invests and causes the firm to swtich its investment from dirty to clean, then the firm will operate at a larger scale than it would have if it were operating clean technology and was financed by return maximizing investors alone. Furthermore, the existence of return-maximizing funds creates a threat to the SR fund, which causes the firm's scale to increase further. The authors call this mechanism complementarity.

\section{Contribution}

The most interesting result of the paper is that even if return-maximizing investors have deep pockets and social investment is costly, SR funds can still have a positive impact through encouraging reform among firms that use dirty technology. This result, however, predates this paper and was first shown by \citet{chowdhry_davies_waters_2019}.

Therefore, the main contribution of this paper is to bring an optimal contracting perspective to the problem of social investment. This leads to two novelties. First the paper finds contracts to implement a switch to green technology under the impact mandate. They propose using either green bonds with covenants that require firms to invest exclusively in the clean technology, or a dual-class share structure that gives all the voting rights to the SR fund. The SR funds must pay a premium for these governance rights.

The second contribution from studying the optimal contracting problem is to understand the role of financial constraints. This yields the complementarity mechanism described earlier.

Finally the authors introduce a social profitability index (SPI). The index just says SR funds should invest where the ``bang for their buck" is highest. Intuitively, this means investing in firms that are almost willing to invest in green technology but just need a small carrot.


\section{Suggestions}

A common prior about SR funds is that they should invest in socially responsible firms. Like \citet{chowdhry_davies_waters_2019}, this paper challenges this belief. They make a compelling argument that SR funds should try to reform firms that use dirty technology. However, this opposite extreme is not realistic either. In this model, there is only a carrot and no stick. Furthermore, the carrot is only given to bad firms. In practice SR funds face a complicated tradeoff between giving the carrot to bad firms, giving the stick to average firms, and help the best firms expand. The very bare bones approach of the SPI does not provide much insight into how SR funds should split money between dirty firms, so instead the paper should focus on how SR funds split their investments between sticks and carrots and between bad firms and good firms.

To create this richer tradeoff, I suggest creating a third type of capital that has a per-unit social benefit. It should also have the highest per-unit investment cost. Call this type of capital ``enhancing capital." Using the example of climate change, enhancing capital could be a machine that creates lithium batteries. If the battery company is constrained, increasing production of batteries could help other companies abate more quickly. This per-unit social benefit captures this effect in reduced form.

I hypothesize the effect of adding enhancing capital would generate the three possible investment strategies described above. In particular an SR fund could:

\begin{itemize}
    \item Give a carrot to a dirty firm
    \item Give a stick to a clean firm for not becoming enhancing
    \item Give a carrot to an enhancing firm
\end{itemize}

The paper already establishes the role of giving a carrot to a dirty firm to become a clean firm. If an enhancing firm is constrained, it is straightforward to see that a carrot to the enhancing firm could be welfare improving. 

Giving a stick to a clean firm is less straightforward. Since this is just a one-period model, there are no dynamics that allow threatening to revoke capital from a clean firm. Instead, just think of giving a stick to a clean firm for not becoming enhancing is like incentivizing a naturally dirty firm to become enhancing in one fell swoop. The stick is disincentivizing the firm from merely becoming clean.

With these three possibilities for an SR fund, the fund's decision problem becomes more interesting and realistic. Although an important part of an SR fund's role is to incentivize dirty firms to become clean, it is also important for them to help scale up constrained enhancing firms. Since the paper does not have enhancing capital, the results make it seem as though SR funds should avoid investing in the most socially beneficial (or least socially costly) firms.

Another implication of having three possibilities is that the narrow SR fund can change the allocation by relaxing the financial constraint of the enhancing firm. Again this feels more realistic. If government spending can improve the allocation, then it makes sense that an SR fund can also improve the allocation.


\newpage
\begin{thebibliography}{9}

\bibitem[Chowdhry et~al.(2019)Chowdhry, Davies, and Waters]{chowdhry_davies_waters_2019}
Chowdhry, Bhagwan, Shaun William Davies, and Brian Waters. 2019.
``Investing for Impact.''
\textit{The Review of Financial Studies} 32 (3): 864--904.
DOI: 10.1093/rfs/hhy068.


\bibitem[Holmstr\"om and Tirole(1997)]{holmstrom_tirole_1997}
Holmstr\"om, Bengt, and Jean Tirole. 1997.
``Financial Intermediation, Loanable Funds, and the Real Sector.''
\textit{The Quarterly Journal of Economics} 112 (3): 663--691.
DOI: 10.1162/003355397555316.

\bibitem[Oehmke and Opp(2025)]{oehmke_opp_2025}
Oehmke, Martin, and Marcus M. Opp. 2025.
``A Theory of Socially Responsible Investment.''
\textit{The Review of Economic Studies} 92 (2): 1193--1225.
DOI: 10.1093/restud/rdae048.

\end{thebibliography}

\end{document}
